
\textbf{لم چگالی اعداد گویا:}
در هر بازهٔ باز حقیقی $(a,b)$ با شرط $a<b$، دست‌کم یک عدد گویا وجود دارد.

چون $a<b$، داریم:
\[
b-a>0
\]

عدد طبیعی $n$ای وجود دارد به‌طوری که:
\[
\frac1n<b-a
\]

پس:
\[
n(b-a)>1
\]

بنابراین:
\[
nb-na>1
\]

پس بین دو عدد حقیقی $na$ و $nb$ دست‌کم یک عدد صحیح وجود دارد؛
یعنی عدد صحیحی مانند $m$ هست که:
\[
na<m<nb
\]

اکنون طرفین را بر $n$ تقسیم می‌کنیم:
\[
a<\frac{m}{n}<b
\]

ولی $\frac{m}{n}$ عددی گویاست، پس در بازهٔ $(a,b)$ یک عدد گویا وجود دارد.

\subsectionaddtolist{الف}

فرض کنید
\(
\mathcal{I}=\{(a_i,b_i)\}_{i\in I}
\)
خانواده‌ای از بازه‌های باز باشد به‌طوری که هیچ دو بازه‌ای اشتراک نداشته باشند.

طبق لم چگالی، هر بازهٔ $(a_i,b_i)$ شامل دست‌کم یک عدد گویاست.  
پس برای هر $i$ عددی گویا مانند
\[
q_i\in(a_i,b_i)\cap\mathbb{Q}
\]
انتخاب می‌کنیم.

اگر $i\neq j$، چون
\[
(a_i,b_i)\cap(a_j,b_j)=\varnothing
\]
پس نمی‌توان یک عدد گویا در هر دو بازه باشد؛ بنابراین
\[
q_i\neq q_j
\]

در نتیجه تابع
\[
i\mapsto q_i
\]
یک‌به‌یک از مجموعهٔ بازه‌ها به $\mathbb{Q}$ است.

از آنجا که $\mathbb{Q}$ شماراست، نتیجه می‌شود مجموعهٔ بازه‌ها نیز شماراست.

\subsectionaddtolist{ب}

می‌دانیم مجموعهٔ
\(
\mathbb{Q}^2=\mathbb{Q}\times\mathbb{Q}
\)
شماراست و در $\mathbb{R}^2$ چگال است.
پس هر دایرهٔ باز شامل نقطه‌ای با مختصات گویاست.
برای هر دایرهٔ $D_i$ نقطه‌ای انتخاب می‌کنیم:
\[
p_i=(r_i,s_i)\in D_i\cap\mathbb{Q}^2
\]

اگر $i\neq j$، چون:
\[
D_i\cap D_j=\varnothing
\]
پس:
\[
p_i\neq p_j
\]

بنابراین نگاشت
\[
D_i\mapsto p_i
\]
یک‌به‌یک به $\mathbb{Q}^2$ است.

چون $\mathbb{Q}^2$ شماراست، نتیجه می‌شود تعداد دایره‌ها نیز شماراست.

\subsectionaddtolist{ج}


فرض کنید $\{D_i\}_{i\in I}$ خانواده‌ای از دایره‌های باز با شعاع مثبت در صفحه باشد به‌طوری که هیچ سه دایره‌ای اشتراک مشترک نداشته باشند.

می‌خواهیم نشان دهیم مجموعهٔ اندیس‌ها $I$ شماراست.

می‌دانیم
\[
\mathbb{Q}^2=\mathbb{Q}\times\mathbb{Q}
\]
شماراست و در $\mathbb{R}^2$ چگال است؛ بنابراین هر دایرهٔ باز شامل دست‌کم یک نقطه با مختصات گویاست.

برای هر دایرهٔ $D_i$ یک نقطهٔ گویا انتخاب می‌کنیم:
\[
p_i\in D_i\cap \mathbb{Q}^2
\]

اکنون برای هر نقطهٔ گویا
\[
p\in\mathbb{Q}^2
\]
مجموعهٔ
\[
A_p=\{\,i\in I \mid p_i=p\,\}
\]
را در نظر بگیرید.

ادعا می‌کنیم:
\[
|A_p|\le 2
\]

زیرا اگر
\[
|A_p|\ge 3،
\]
آنگاه سه اندیس متمایز
\[
i_1,i_2,i_3
\]
وجود دارند به‌طوری که
\[
p_{i_1}=p_{i_2}=p_{i_3}=p.
\]

طبق تعریف نقاط انتخابی داریم
\[
p\in D_{i_1}\cap D_{i_2}\cap D_{i_3},
\]
پس سه دایرهٔ
\[
D_{i_1},D_{i_2},D_{i_3}
\]
اشتراک مشترک دارند که با فرض مسئله تناقض است.

پس هر نقطهٔ گویا حداکثر متناظر با دو دایره است.

اکنون چون $\mathbb{Q}^2$ شماراست، می‌توان نوشت
\[
\mathbb{Q}^2=\{p_1,p_2,p_3,\dots\}.
\]

در نتیجه:
\[
I=\bigcup_{n=1}^{\infty} A_{p_n},
\]
و هر مجموعهٔ $A_{p_n}$ متناهی (در واقع با حداکثر دو عضو) است.

اما اجتماع شمارای مجموعه‌های متناهی، شماراست؛ بنابراین $I$ شماراست.

پس تعداد دایره‌ها شماراست.

